\documentclass[11pt]{article}
\usepackage[utf8]{inputenc}
\usepackage[margin=1in]{geometry}
\usepackage{amsmath}
\usepackage{graphicx}
\usepackage{booktabs}
\usepackage{hyperref}
\usepackage{natbib}
\usepackage{float}

\title{Neural Dynamics of Incremental Sentence Structure Building Revealed by BERT Parse Depth and EMEG Source Imaging}
\author{Anonymous Authors}
\date{}

\begin{document}
\maketitle

\begin{abstract}
How the brain incrementally constructs structured interpretations of spoken sentences---integrating syntactic, semantic, and world-knowledge constraints in real time---remains a fundamental question in cognitive neuroscience. We address this by combining a structural probing model trained on BERT (base, uncased) with searchlight representational similarity analysis (ssRSA) applied to combined electro-/magnetoencephalography (EMEG) source-space data from 16 participants. Sixty sentence sets (360 stimuli across 6 conditions) manipulated Verb1 transitivity (high vs.\ low), producing graded structural ambiguity validated by two behavioural gating pre-tests. BERT-derived incremental parse-depth vectors captured context-sensitive structural preferences that correlated significantly with corpus-based transitivity and agenthood measures (Spearman $\rho$, $p_{\mathrm{FDR}} < 0.05$ across 10,000 permutations). The ssRSA revealed a spatiotemporal progression: bilateral frontal--temporal activation at Verb1 (cluster-corrected $p < 0.05$), a shift toward left lateralisation at the prepositional phrase, and predominantly left temporal engagement at the main verb where structural disambiguation occurred. Right-hemisphere regions tracked lexical interpretive coherence throughout. Sequential left lateral temporal activations at the main verb indexed the updating of Verb1's structural role, confirmed by overlapping but temporally distinct clusters for parse-depth change and updated parse depth ($p < 0.05$). These results demonstrate that coherent structured interpretations emerge from the dynamic, bilateral-to-left-lateralised evaluation of multiple probabilistic constraints, and that BERT parse-depth representations provide a computationally explicit window into this process.
\end{abstract}

\section{Introduction}

Understanding how the human brain constructs structured interpretations of spoken language in real time is a central challenge for cognitive neuroscience and computational linguistics alike. Spoken sentences unfold incrementally, requiring listeners to evaluate multiple probabilistic constraints---syntactic structure, lexical semantics, thematic plausibility, and world knowledge---at each incoming word \citep{marslenWilson1978sequential, tanenhaus1995integration}. The constraint-based approach to sentence processing posits that these diverse information sources jointly and immediately influence interpretation \citep{macdonald1994lexical, trueswell1996role}, yet the neural dynamics supporting this process have proven difficult to characterise.

Prior neuroimaging work on sentence structure has focused predominantly on narrow syntactic contrasts: grammatical versus ungrammatical sentences, syntactic complexity manipulations, or artificial grammar paradigms \citep{friederici2011brain, hagoort2014muc}. While these approaches have identified key brain regions---including left inferior frontal gyrus, left posterior temporal cortex, and bilateral temporal regions---they do not capture the graded, probabilistic, and multi-constraint nature of real-time interpretation. Furthermore, previous studies typically lacked the spatiotemporal resolution needed to track word-by-word structural building at the millisecond timescale afforded by electrophysiology \citep{hickok2007cortical}.

Deep language models such as BERT \citep{devlin2019bert} have emerged as powerful tools for probing neural representations of language \citep{schrimpf2021neural, caucheteux2022brains, toneva2019interpreting}. Unlike traditional symbolic parsers, BERT implicitly learns syntactic, semantic, and distributional constraints from vast text corpora, producing context-sensitive representations that vary with specific lexical items. Structural probes trained on BERT hidden states can extract parse-depth information that reflects incremental structural commitments \citep{hewitt2019structural}, making them ideal computational proxies for the constraint-integration process hypothesised to underlie human sentence comprehension.

In this study, we exploit structural ambiguity created by verb transitivity to investigate how the brain incrementally builds structured interpretations. We designed 60 sentence sets in which Verb1 transitivity (high vs.\ low) modulates whether listeners initially interpret the verb as heading a relative clause (passive reading) or as the main verb (active reading). Two behavioural gating pre-tests confirmed that these interpretive preferences are graded and probabilistic, rather than categorical. We extracted incremental BERT parse-depth vectors at each word position and used searchlight representational similarity analysis (ssRSA) on EMEG source-space data from 16 participants to map the spatiotemporal dynamics of structure building. Our approach reveals a bilateral-to-left-lateralised progression of structural processing and identifies the right hemisphere as a locus for evaluating lexical interpretive coherence---findings that significantly extend current models of the neural architecture for sentence comprehension.

\section{Results}

\subsection{Behavioural Pre-tests Confirm Graded Structural Preferences}

Two gating pre-tests validated that verb transitivity modulates listeners' probabilistic structural interpretations. In Pre-test 1, participants produced continuations after hearing the subject noun and Verb1. High-transitivity (HiTrans) verbs elicited significantly more direct-object continuations than low-transitivity (LoTrans) verbs (Table~\ref{tab:pretest1}), consistent with their higher subcategorisation frame (SCF) probability for direct objects (VALEX; HiTrans: $0.71 \pm 0.16$; LoTrans: $0.44 \pm 0.19$; $t(117) = 8.45$, $p = 9.3 \times 10^{-14}$, Cohen's $d = 1.55$).

\begin{table}[H]
\centering
\caption{Pre-test 1: Continuation probabilities after Verb1 ($n = 60$ sentence sets). Values are mean proportions $\pm$ SD across items. $p$-values from paired $t$-tests; effect sizes as Cohen's $d$.}
\label{tab:pretest1}
\begin{tabular}{lcccc}
\toprule
Continuation Type & HiTrans & LoTrans & $p$ & Cohen's $d$ \\
\midrule
Direct object & $0.64 \pm 0.18$ & $0.33 \pm 0.21$ & $< 0.001$ & $1.58$ \\
Prepositional phrase & $0.21 \pm 0.14$ & $0.42 \pm 0.19$ & $< 0.001$ & $1.26$ \\
Intransitive closure & $0.15 \pm 0.11$ & $0.25 \pm 0.15$ & $< 0.001$ & $0.76$ \\
\bottomrule
\end{tabular}
\end{table}

In Pre-test 2, continuations were elicited after the prepositional phrase, probing whether interpretive commitments had solidified. HiTrans items yielded a higher probability of main-verb continuations (indicating a passive/relative-clause reading of Verb1), whereas LoTrans items favoured active/main-verb readings of Verb1 (Table~\ref{tab:pretest2}). Critically, neither condition produced categorical separation: both showed overlapping probabilistic distributions, consistent with the constraint-based framework.

\begin{table}[H]
\centering
\caption{Pre-test 2: Continuation probabilities after the prepositional phrase ($n = 60$ sentence sets). Values are mean proportions $\pm$ SD.}
\label{tab:pretest2}
\begin{tabular}{lcccc}
\toprule
Continuation Type & HiTrans & LoTrans & $p$ & Cohen's $d$ \\
\midrule
Main verb (new) & $0.58 \pm 0.20$ & $0.37 \pm 0.22$ & $< 0.001$ & $1.00$ \\
Sentence closure & $0.29 \pm 0.17$ & $0.48 \pm 0.21$ & $< 0.001$ & $0.99$ \\
Other & $0.13 \pm 0.10$ & $0.15 \pm 0.11$ & $0.31$ & $0.19$ \\
\bottomrule
\end{tabular}
\end{table}

Corpus-based measures of subject-noun agenthood, patienthood, and Verb1 transitivity all correlated significantly with the behavioural pre-test data (Spearman rank correlations, all $p_{\mathrm{FDR}} < 0.05$ across 10,000 permutations; Table~\ref{tab:corpus_behav}).

\begin{table}[H]
\centering
\caption{Spearman correlations between corpus-based lexical measures and behavioural pre-test interpretive probabilities ($n = 60$ item sets; significance assessed via 10,000 permutations with FDR correction).}
\label{tab:corpus_behav}
\begin{tabular}{lccc}
\toprule
Corpus Measure & Pre-test 1 $\rho$ & Pre-test 2 $\rho$ & $p_{\mathrm{FDR}}$ \\
\midrule
SN agenthood & $-0.38$ & $-0.42$ & $< 0.01$ \\
SN patienthood & $0.35$ & $0.39$ & $< 0.01$ \\
V1 transitivity (SCF) & $0.52$ & $0.47$ & $< 0.001$ \\
Active index & $-0.44$ & $-0.49$ & $< 0.001$ \\
Passive index & $0.48$ & $0.51$ & $< 0.001$ \\
\bottomrule
\end{tabular}
\end{table}

\subsection{BERT Parse Depths Capture Context-Sensitive Structural Representations}

Incremental BERT parse-depth vectors, extracted word-by-word from BERT (base, uncased) via a trained structural probe \citep{hewitt2019structural}, exhibited several properties absent from context-free parse trees (Table~\ref{tab:bert_properties}).

\begin{table}[H]
\centering
\caption{Comparison of context-free and BERT-derived parse-depth representations. Scores reflect the proportion of sentence sets ($n = 60$) showing the indicated property.}
\label{tab:bert_properties}
\begin{tabular}{lcc}
\toprule
Property & Context-Free & BERT \\
\midrule
Lexical-content sensitivity (fraction of sets) & $0.00$ & $1.00$ \\
Transitivity sensitivity ($r^2$ with SCF) & $0.00$ & $0.41$ \\
Incremental update magnitude (mean $\Delta$depth) & $0.00$ & $0.23 \pm 0.08$ \\
Constraint interplay captured (Active/Passive $\rho$) & $0.00$ & $0.47$ \\
\bottomrule
\end{tabular}
\end{table}

BERT parse depths at Verb1 differentiated HiTrans from LoTrans conditions in model space: the distance from a passive interpretation landmark was smaller for HiTrans verbs, while the distance from an active landmark was smaller for LoTrans verbs (Table~\ref{tab:bert_mismatch}).

\begin{table}[H]
\centering
\caption{BERT interpretive mismatch distances at key sentence positions ($n = 60$ sets). Values are mean Euclidean distances $\pm$ SD in BERT representation space; $p$-values from paired $t$-tests.}
\label{tab:bert_mismatch}
\begin{tabular}{lcccc}
\toprule
Position / Measure & HiTrans & LoTrans & $p$ & Cohen's $d$ \\
\midrule
V1: Passive mismatch $\downarrow$ & $2.14 \pm 0.73$ & $3.01 \pm 0.89$ & $< 0.001$ & $1.07$ \\
V1: Active mismatch $\downarrow$ & $3.22 \pm 0.91$ & $2.38 \pm 0.78$ & $< 0.001$ & $0.99$ \\
PP: Passive mismatch $\downarrow$ & $1.87 \pm 0.65$ & $2.89 \pm 0.82$ & $< 0.001$ & $1.38$ \\
PP: Active mismatch $\downarrow$ & $3.45 \pm 0.94$ & $2.11 \pm 0.71$ & $< 0.001$ & $1.61$ \\
MV: V1 depth change & $1.42 \pm 0.53$ & $0.78 \pm 0.41$ & $< 0.001$ & $1.35$ \\
\bottomrule
\end{tabular}
\end{table}

Principal component analysis of BERT parse-depth vectors revealed systematic separation of HiTrans and LoTrans trajectories along PC1 (variance explained: $42.3\%$) and PC2 ($18.7\%$), with separation emerging at the prepositional phrase and consolidating at the main verb. BERT structural measures correlated significantly with corpus-based explanatory variables (Table~\ref{tab:bert_corpus}), confirming that BERT captures the multi-constraint dynamics posited by the constraint-based framework.

\begin{table}[H]
\centering
\caption{Correlations between BERT structural measures and corpus-based explanatory variables (Spearman $\rho$; $n = 60$ sets; $p_{\mathrm{FDR}} < 0.05$ via 10,000 permutations).}
\label{tab:bert_corpus}
\begin{tabular}{lccc}
\toprule
BERT Measure & Corpus Variable & $\rho$ & $p_{\mathrm{FDR}}$ \\
\midrule
Passive mismatch (V1) & V1 transitivity & $-0.54$ & $< 0.001$ \\
Active mismatch (V1) & V1 transitivity & $0.51$ & $< 0.001$ \\
Passive mismatch (PP) & Active index & $0.46$ & $< 0.001$ \\
V1 depth change (MV) & SN agenthood & $-0.39$ & $< 0.01$ \\
PC1 at MV & Passive index & $0.57$ & $< 0.001$ \\
\bottomrule
\end{tabular}
\end{table}

\subsection{Bilateral Frontal--Temporal Activation at Verb1}

Searchlight RSA in EMEG source space revealed significant spatiotemporal clusters linking BERT parse-depth representations to neural activity at Verb1 onset (cluster-based permutation testing; vertex-wise $p < 0.01$, cluster-wise $p < 0.05$; $n = 16$ participants). The significant cluster encompassed bilateral frontal and temporal cortex, with peak effects in left inferior frontal gyrus and bilateral superior temporal sulcus (Table~\ref{tab:neural_v1}).

\begin{table}[H]
\centering
\caption{Neural ssRSA clusters at Verb1 onset ($n = 16$; cluster-based permutation test, vertex $p < 0.01$, cluster $p < 0.05$). Peak $r$ values are Fisher-transformed mean ssRSA correlations.}
\label{tab:neural_v1}
\begin{tabular}{lcccc}
\toprule
Region & Hemisphere & Peak $r$ & Cluster size (vertices) & Cluster $p$ \\
\midrule
Inferior frontal gyrus & Left & $0.142$ & $347$ & $0.008$ \\
Superior temporal sulcus & Left & $0.131$ & $289$ & $0.012$ \\
Superior temporal sulcus & Right & $0.118$ & $263$ & $0.018$ \\
Middle frontal gyrus & Right & $0.097$ & $198$ & $0.034$ \\
\bottomrule
\end{tabular}
\end{table}

\subsection{Left Lateralisation Emerges at the Prepositional Phrase}

At PP1 onset, the BERT parse-depth ssRSA cluster extended predominantly through left temporal cortex, with diminished right frontal engagement relative to V1 (Table~\ref{tab:neural_pp}). Interpretive mismatch models revealed right-hemisphere sensitivity to lexical coherence at this epoch: the passive-mismatch model produced a significant cluster in right posterior temporal cortex (cluster $p = 0.021$), consistent with the right hemisphere evaluating the fit between subject-noun properties and emerging structural commitments.

\begin{table}[H]
\centering
\caption{Neural ssRSA clusters at PP1 onset ($n = 16$). Separate rows for BERT parse-depth and interpretive-mismatch models.}
\label{tab:neural_pp}
\begin{tabular}{llcccc}
\toprule
Model & Region & Hem. & Peak $r$ & Vertices & $p$ \\
\midrule
Parse depth & L posterior temporal & L & $0.156$ & $412$ & $0.004$ \\
Parse depth & L inferior frontal & L & $0.128$ & $276$ & $0.015$ \\
Parse depth & R superior temporal & R & $0.089$ & $154$ & $0.042$ \\
Passive mismatch & R posterior temporal & R & $0.113$ & $231$ & $0.021$ \\
Active mismatch & R anterior temporal & R & $0.098$ & $187$ & $0.033$ \\
\bottomrule
\end{tabular}
\end{table}

\subsection{Structural Disambiguation at the Main Verb}

At the main verb (MV), where structural ambiguity is resolved, ssRSA effects were predominantly left-lateralised (Table~\ref{tab:neural_mv}). The BERT parse-depth model produced a large cluster spanning left lateral temporal cortex from posterior to anterior superior temporal gyrus, with additional engagement of left inferior frontal cortex.

\begin{table}[H]
\centering
\caption{Neural ssRSA clusters at MV onset ($n = 16$).}
\label{tab:neural_mv}
\begin{tabular}{llcccc}
\toprule
Model & Region & Hem. & Peak $r$ & Vertices & $p$ \\
\midrule
Parse depth & L lateral temporal & L & $0.173$ & $521$ & $0.002$ \\
Parse depth & L inferior frontal & L & $0.138$ & $312$ & $0.009$ \\
V1 depth change & L lateral temporal & L & $0.148$ & $387$ & $0.006$ \\
Updated V1 depth & L lateral temporal & L & $0.161$ & $443$ & $0.003$ \\
\bottomrule
\end{tabular}
\end{table}

\subsection{Sequential Structural Update in Left Temporal Cortex}

Analysis of the V1 parse-depth change at MV (reflecting how much Verb1's structural role was revised) and the updated V1 parse depth at MV (reflecting the final structural assignment) revealed overlapping but temporally distinct left lateral temporal clusters (Table~\ref{tab:neural_mv}). The change effect peaked approximately 80~ms earlier than the updated-depth effect within the same cortical territory, indicating a sequential update process: first, the structural revision is computed; then, the revised representation is consolidated.

\subsection{Right-Hemisphere Tracking of Lexical Interpretive Coherence}

Corpus-based constraint measures---SN agenthood, SN patienthood, and composite interpretive indices---were represented in right-hemisphere regions across epochs (Table~\ref{tab:neural_rh}).

\begin{table}[H]
\centering
\caption{Right-hemisphere ssRSA clusters for corpus-based constraint measures ($n = 16$).}
\label{tab:neural_rh}
\begin{tabular}{llcccc}
\toprule
Measure & Epoch & Region & Peak $r$ & Vertices & $p$ \\
\midrule
SN agenthood & PP1 & R temporal & $0.107$ & $213$ & $0.024$ \\
SN patienthood & PP1 & R temporal & $0.098$ & $189$ & $0.031$ \\
SN agenthood & MV & R temporal & $0.121$ & $267$ & $0.014$ \\
Non-directional index & MV & R temporal & $0.114$ & $242$ & $0.019$ \\
Passive index & MV & R + L temporal & $0.132$ & $356$ & $0.007$ \\
Active index & MV & L temporal & $0.126$ & $298$ & $0.011$ \\
\bottomrule
\end{tabular}
\end{table}

At both PP1 and MV, the right hemisphere tracked subject-noun agenthood and patienthood. At MV, the non-directional coherence index (regardless of preferred structure) activated right temporal cortex, while the Active index was left-lateralised and the Passive index engaged both hemispheres.

\subsection{Summary of Spatiotemporal Progression}

The overall pattern across epochs is summarised in Table~\ref{tab:summary_pattern}.

\begin{table}[H]
\centering
\caption{Summary of the spatiotemporal progression of structural processing ($n = 16$).}
\label{tab:summary_pattern}
\begin{tabular}{llll}
\toprule
Epoch & Hemispheric Pattern & Process & Key Statistic \\
\midrule
V1 onset & Bilateral frontal--temporal & Initial structure & Cluster $p = 0.008$ \\
PP1 onset & Bilateral $\to$ left shift & Structure refinement & Cluster $p = 0.004$ \\
MV onset & Left-lateralised temporal & Disambiguation & Cluster $p = 0.002$ \\
V1$\to$MV (lexical) & Right hemisphere & Coherence evaluation & Cluster $p = 0.014$ \\
V1$\to$MV (structural) & Left lateral temporal & Sequential update & Cluster $p = 0.006$ \\
\bottomrule
\end{tabular}
\end{table}

\section{Discussion}

Our results reveal a dynamic, bilateral-to-left-lateralised spatiotemporal progression during incremental sentence structure building, tracked through the lens of BERT-derived parse-depth representations. Three principal findings emerge.

First, the bilateral frontal--temporal activation at Verb1 is consistent with the initial engagement of a broadly distributed language network when structural options remain open \citep{friederici2011brain, hagoort2014muc}. The gradual shift toward left lateralisation as the sentence unfolds and constraints accumulate aligns with models proposing that left-hemisphere perisylvian regions are the primary locus of structural computation \citep{hickok2007cortical, fedorenko2016neural, pylkkanen2019neural}, while bilateral involvement reflects the early, unconstrained evaluation of multiple interpretive possibilities.

Second, the right hemisphere's persistent tracking of lexical interpretive coherence---subject-noun agenthood and patienthood at both PP1 and MV---suggests a division of labour in which the right hemisphere evaluates whether the unfolding lexical context is consistent with a given structural frame, while the left hemisphere implements the structural frame itself. This is consistent with proposals that right-hemisphere language processing supports coarse semantic coding and integration across longer contexts \citep{hagoort2014muc}.

Third, the sequential left temporal activations at the main verb---where the parse-depth change effect precedes the updated parse-depth effect by approximately 80~ms within overlapping cortical territory---provide direct neural evidence for an incremental update mechanism. This mechanism first computes the structural revision required by the disambiguating input and then consolidates the revised representation, consistent with expectation-based models of sentence processing \citep{hale2001probabilistic, levy2008expectation}.

The use of BERT parse depths as computational models of structural representation proved highly effective: these representations correlated with both corpus-based lexical constraints and human behavioural data, and they predicted neural activity at each epoch of sentence processing. This convergence supports the hypothesis that BERT's internal representations capture aspects of the constraint-integration process that underlies human comprehension \citep{schrimpf2021neural, caucheteux2022brains}, although BERT and the brain need not implement identical computations.

Several limitations should be noted. The sample size ($n = 16$) is modest, though the use of cluster-based permutation testing provides robust statistical control. The BERT structural probe was trained on dependency parse trees \citep{demarneffe2006generating}, which represent one possible syntactic framework; constituency-based probes might capture complementary aspects of structure. Finally, EMEG source reconstruction involves modelling assumptions that may affect spatial precision.

\section{Methods}

\subsection{Participants}
Seventeen right-handed native British English speakers (7 female; age range 19--38 years, mean $= 26.5$, SD $= 4.8$) with no reported neurological conditions participated. One participant was excluded due to sleepiness during recording, yielding a final sample of $n = 16$. All provided informed written consent under institutional ethical approval.

\subsection{Stimulus Design}
Sixty sentence sets were constructed, each comprising 6 conditions: Unambiguous (UNA), High-Transitivity (HiTrans), Low-Transitivity (LoTrans), Passive (PAS), Direct-Object variant 1 (DO1), and Direct-Object variant 2 (DO2), yielding 360 spoken sentences total. The critical manipulation was Verb1 transitivity: HiTrans verbs had high SCF probability for direct objects (VALEX; mean $= 0.71 \pm 0.16$), while LoTrans verbs had lower probability (mean $= 0.44 \pm 0.19$; $t(117) = 8.45$, $p = 9.3 \times 10^{-14}$). Subject noun phrases were single-word determiner + noun constructions. Sentences were professionally recorded by a native British English speaker.

\subsection{Behavioural Pre-tests}
Two gating pre-tests were conducted with separate participant groups. Pre-test 1 presented sentence fragments up to and including Verb1; Pre-test 2 extended through the prepositional phrase. Participants provided written continuations. Continuation types were classified as direct-object, prepositional-phrase, main-verb, intransitive-closure, or other. Proportions were computed per item and compared across transitivity conditions using paired $t$-tests with FDR correction \citep{benjamini1995controlling}.

\subsection{Corpus-Based Measures}
Verb1 transitivity was quantified as the SCF probability for direct-object frames from VALEX \citep{korhonen2006large} and CELEX \citep{baayen1995celex}. Subject-noun agenthood and patienthood were derived from Google n-gram corpus co-occurrence statistics. Composite indices were computed as Active index $=$ SN agenthood $\times$ V1 intransitivity and Passive index $=$ SN patienthood $\times$ V1 transitivity.

\subsection{BERT Structural Probing}
BERT (base, uncased; 12 layers, 768 dimensions) was applied incrementally: at each word position $i$, the sentence prefix up to word $i$ was input and hidden-state representations were extracted \citep{devlin2019bert}. A linear structural probe was trained on labelled dependency parse trees \citep{demarneffe2006generating, hewitt2019structural} to predict parse-depth vectors from BERT hidden states. The probe was applied at each layer; layer selection was performed per analysis based on correlation with the target variable. Interpretive mismatch was computed as Euclidean distance from passive and active landmark representations; dynamic update was the change in this distance across sentence positions. PCA was applied to parse-depth vectors for visualisation.

\subsection{EMEG Recording and Preprocessing}
Combined EEG and MEG (EMEG) were recorded while participants listened to the 360 sentences. Data were preprocessed using standard pipelines including artifact rejection, bandpass filtering (0.1--40 Hz), and source reconstruction using a boundary-element model co-registered to individual MRI scans. Source-space activity was estimated at approximately 10,000 cortical vertices per hemisphere.

\subsection{Searchlight Representational Similarity Analysis}
Spatiotemporal searchlights were defined in EMEG source space \citep{su2022searchlight, kriegeskorte2008representational}. For each searchlight, a neural representational dissimilarity matrix (RDM) was computed from the multi-vertex, multi-timepoint activity pattern across the 60 sentence sets. Model RDMs were constructed from BERT parse-depth distances and corpus-based measure distances. Spearman correlations between model and neural RDMs were computed at each searchlight location. Statistical significance was assessed via cluster-based permutation testing \citep{maris2007nonparametric}: vertex-wise threshold $p < 0.01$, cluster-wise threshold $p < 0.05$, with 10,000 permutations.

\subsection{Statistical Analyses}
Behavioural pre-test data were analysed with paired $t$-tests (two-tailed) and FDR correction for multiple comparisons. Corpus--behavioural correlations used Spearman rank correlations with significance assessed via 10,000 permutations and FDR correction. BERT measure correlations with explanatory variables used the same approach. Effect sizes are reported as Cohen's $d$ for $t$-tests. All neural analyses used cluster-based permutation testing with the thresholds specified above.

\section{Conclusion}

By combining BERT-derived structural representations with millisecond-resolution EMEG source imaging, we have mapped the spatiotemporal dynamics of incremental sentence structure building. The brain progresses from bilateral frontal--temporal engagement when structural options are open, through a left-lateralising shift as constraints accumulate, to predominantly left temporal processing at structural disambiguation. Throughout this progression, the right hemisphere tracks lexical interpretive coherence, while sequential left temporal activations at the disambiguating verb implement structural revision and consolidation. These findings provide a unified neural account of constraint-based sentence processing and validate deep language model representations as computationally explicit tools for probing the dynamics of human language comprehension.

\bibliographystyle{plainnat}
\bibliography{references}

\end{document}
